% GENERATED FILE. Edit canonical lessons, then run npm run generate:books.
% canonical-source-hash: fnv1a64:4849b11f0560081a
% canonical-lessons: ZH-C191-huilai, ZH-W191-hui-return, ZH-C191-huiqu, ZH-C191-huiguo, ZH-C191-chuguo, ZH-C191-chuchai

\chapter{Back Next Month}
\label{ch:zh-back-next-month}
\hlchaptermodality{191}

\begin{chapteropening}
\textbf{By the end of this chapter:} \emph{I can say when I will come back or go back, go home from abroad, go abroad or go away for work.}

Complete the last lesson of chapter 191: I can say when I will come back or go back, go home from abroad, go abroad or go away for work.
\end{chapteropening}

\section[\texorpdfstring{\zh{回来} (huílái)}{回来 (huílái)}]{\zh{回来} (huílái) — to come back}

\label{lesson:ZH-C191-huilai}



\begin{quote}
Before the new one: say the Mandarin for cinema, then the Mandarin for museum.
\end{quote}

\subsection*{You'll want to know: \zh{回来}}
\textbf{\zh{回来}} — \emph{huílái} — \textquotedblleft{}to come back\textquotedblright{}.

\begin{quote}
\textbf{\zh{我下个月五号回来}} \emph{wǒ xià ge yuè wǔ hào huílái} — I'll be back on the fifth of next month.
\end{quote}

The new character, in pieces: \zh{回} — six strokes: a large frame, a small \zh{口} inside it, and the bar that closes the large frame last. The next lesson puts it in your hand.

\begin{cousinweb}
\textbf{\zh{回来}} (\emph{huílái}) is built from \zh{回} 'to turn back' and \zh{来} 'to come', both inherited from Middle Chinese, \emph{hwoj} and \emph{loj} (Guangyun, in Baxter's transcription), which glosses the first 'to turn round'.
\end{cousinweb}

\begin{grammarlens}[title={What you've built}]
One more word for talking about plans.
\end{grammarlens}

\subsection*{Your turn}
\begin{itemize}
\raggedright
  \item \emph{Say it:} \emph{huílái}
  \item \emph{Say it:} \emph{huílái}, once more
  \item \emph{Say it:} \emph{wǒ xià ge yuè wǔ hào huílái} — I'll be back on the fifth of next month.
  \item \emph{Recall:} say the Mandarin for cinema, then the Mandarin for museum, then say \emph{huílái} again
\end{itemize}

\begin{tcolorbox}[breakable,colback=teal!4,colframe=teal!35!black,title={Before you move on}]
What does \zh{回来} mean? (\textquotedblleft{}To come back\textquotedblright{}.) Say it once more.
\end{tcolorbox}
\section[\texorpdfstring{\zh{回} (huí)}{回 (huí)}]{\zh{回} — \zh{口} inside a larger frame}

\label{lesson:ZH-W191-hui-return}



\begin{quote}
The word first, then the character inside it.

\begin{itemize}
\raggedright
  \item \emph{Recall:} say \emph{huílái} — \textbf{R1}, one lesson back, before it decays
\end{itemize}

\begin{quote}
\zh{国} \zh{来}
\end{quote}

\emph{Read it:} each one aloud, with nothing said first — \emph{guó}, \emph{lái}
\end{quote}

\subsection*{Script — \zh{口} inside a larger frame}
\hlblockfigure{\detokenize{figures/ZH-W191-hui-return-filmstrip.pdf}}{How \zh{回} is written, stroke by stroke}

\begin{quote}
\zh{回}
\end{quote}

\textbf{Two pieces: a large frame, and \zh{口} inside it.} Count the pieces, not the strokes.

\textbf{The frame, first two strokes.} A vertical down the left side, then a stroke along the top that turns down the right side.

\textbf{\zh{口} inside, three strokes:} a short vertical, a stroke that runs right and turns down, and the bar that closes it.

\textbf{Last, the frame's foot.} One bar along the bottom closes the large frame, the way the last stroke of \zh{国} does.

2 + 3 + 1 = six.

\subsection*{Your turn}
\begin{itemize}
\raggedright
  \item \emph{Write it:} \zh{回}, all six, with the model in front of you
  \item \emph{Cover:} write it again from memory, and count the pieces, not the strokes
\end{itemize}

\begin{tcolorbox}[breakable,colback=teal!4,colframe=teal!35!black,title={Before you move on}]
How many strokes in \zh{回}? (\textbf{Six}.) What goes inside the frame? (\textbf{\zh{口}}.) Which stroke comes last? (\textbf{The bar that closes the large frame}.)

Source: \href{https://raw.githubusercontent.com/chanind/hanzi-writer-data/68d10a4b21150cae5e1ebbd223eed289cf32d90c/data/\%E5\%9B\%9E.json}{Hanzi Writer Data \zh{回}, pinned stroke paths and medians}.
\end{tcolorbox}
\section[\texorpdfstring{\zh{回去} (huíqù)}{回去 (huíqù)}]{\zh{回去} (huíqù) — to go back there}

\label{lesson:ZH-C191-huiqu}



\begin{quote}
Before the new one: say the Mandarin for museum, then the Mandarin for to come back.
\end{quote}

\subsection*{You'll want to know: \zh{回去}}
\textbf{\zh{回去}} — \emph{huíqù} — \textquotedblleft{}to go back there\textquotedblright{}.

\begin{quote}
\textbf{\zh{我不想回去}} \emph{wǒ bù xiǎng huíqù} — I don't want to go back.
\end{quote}

\begin{cousinweb}
\textbf{\zh{回去}} (\emph{huíqù}) is built from \zh{回} 'to turn back' and \zh{去} 'to go', both inherited from Middle Chinese, \emph{hwoj} and \emph{khjoH} (Guangyun, in Baxter's transcription), which glosses the second 'to leave'.
\end{cousinweb}

\begin{grammarlens}[title={What you've built}]
One more word for talking about plans.
\end{grammarlens}

\subsection*{Your turn}
\begin{itemize}
\raggedright
  \item \emph{Say it:} \emph{huíqù}
  \item \emph{Say it:} \emph{huíqù}, once more
  \item \emph{Say it:} \emph{wǒ bù xiǎng huíqù} — I don't want to go back.
  \item \emph{Recall:} say the Mandarin for museum, then the Mandarin for to come back, then say \emph{huíqù} again
\end{itemize}

\begin{tcolorbox}[breakable,colback=teal!4,colframe=teal!35!black,title={Before you move on}]
What does \zh{回去} mean? (\textquotedblleft{}To go back there\textquotedblright{}.) Say it once more.
\end{tcolorbox}
\section[\texorpdfstring{\zh{回国} (huíguó)}{回国 (huíguó)}]{\zh{回国} (huíguó) — to go home from abroad; to return to one's own country}

\label{lesson:ZH-C191-huiguo}



\begin{quote}
Before the new one: say the Mandarin for to come back, then the Mandarin for to go back there.
\end{quote}

\subsection*{You'll want to know: \zh{回国}}
\textbf{\zh{回国}} — \emph{huíguó} — \textquotedblleft{}to go home from abroad; to return to one's own country\textquotedblright{}.

\begin{quote}
\textbf{\zh{他下星期回国}} \emph{tā xià xīngqī huíguó} — He's going back to his own country next week.
\end{quote}

\begin{cousinweb}
\textbf{\zh{回国}} (\emph{huíguó}) is built from \zh{回} 'to turn back' and \zh{国} 'country', both inherited from Middle Chinese, \emph{hwoj} and \emph{kwok} (Guangyun, in Baxter's transcription).
\end{cousinweb}

\begin{grammarlens}[title={What you've built}]
One more word for talking about plans.
\end{grammarlens}

\subsection*{Your turn}
\begin{itemize}
\raggedright
  \item \emph{Say it:} \emph{huíguó}
  \item \emph{Say it:} \emph{huíguó}, once more
  \item \emph{Say it:} \emph{tā xià xīngqī huíguó} — He's going back to his own country next week.
  \item \emph{Recall:} say the Mandarin for to come back, then the Mandarin for to go back there, then say \emph{huíguó} again
\end{itemize}

\begin{tcolorbox}[breakable,colback=teal!4,colframe=teal!35!black,title={Before you move on}]
What does \zh{回国} mean? (\textquotedblleft{}To go home from abroad; to return to one's own country\textquotedblright{}.) Say it once more.
\end{tcolorbox}
\section[\texorpdfstring{\zh{出国} (chūguó)}{出国 (chūguó)}]{\zh{出国} (chūguó) — to go abroad}

\label{lesson:ZH-C191-chuguo}



\begin{quote}
Before the new one: say the Mandarin for to go back there, then the Mandarin for to go home from abroad.
\end{quote}

\subsection*{You'll want to know: \zh{出国}}
\textbf{\zh{出国}} — \emph{chūguó} — \textquotedblleft{}to go abroad\textquotedblright{}.

\begin{quote}
\textbf{\zh{我明年要出国}} \emph{wǒ míngnián yào chūguó} — I'm going abroad next year.
\end{quote}

\begin{cousinweb}
\textbf{\zh{出国}} (\emph{chūguó}) is built from \zh{出} 'to go out' and \zh{国} 'country', both inherited from Middle Chinese, \emph{tsyhwit} and \emph{kwok} (Guangyun, in Baxter's transcription).
\end{cousinweb}

\begin{grammarlens}[title={What you've built}]
One more word for talking about plans.
\end{grammarlens}

\subsection*{Your turn}
\begin{itemize}
\raggedright
  \item \emph{Say it:} \emph{chūguó}
  \item \emph{Say it:} \emph{chūguó}, once more
  \item \emph{Say it:} \emph{wǒ míngnián yào chūguó} — I'm going abroad next year.
  \item \emph{Recall:} say the Mandarin for to go back there, then the Mandarin for to go home from abroad, then say \emph{chūguó} again
\end{itemize}

\begin{tcolorbox}[breakable,colback=teal!4,colframe=teal!35!black,title={Before you move on}]
What does \zh{出国} mean? (\textquotedblleft{}To go abroad\textquotedblright{}.) Say it once more.
\end{tcolorbox}
\section[chūchāi]{chūchāi — to travel for work; to go on a business trip}

\label{lesson:ZH-C191-chuchai}



\begin{quote}
Before the new one: say the Mandarin for to go home from abroad, then the Mandarin for to go abroad.
\end{quote}

\subsection*{You'll want to know: chūchāi}
\emph{chūchāi} — \textquotedblleft{}to travel for work; to go on a business trip\textquotedblright{}.

\begin{quote}
\emph{bàba hòutiān yào chūchāi} — Dad is going away on business the day after tomorrow.
\end{quote}

Its characters come later; for now it is a word for the ear.

\begin{cousinweb}
\emph{chūchāi} is built from \zh{出} 'to go out' and \emph{chāi} 'an errand'. The first is inherited from Middle Chinese \emph{tsyhwit} (Guangyun, in Baxter's transcription); the dictionary lists the second under several readings, glossed 'uneven' and 'to choose', and gives no errand, so how the word came to mean a work trip is not settled.
\end{cousinweb}

\begin{grammarlens}[title={What you've built}]
One more word for talking about plans.
\end{grammarlens}

\subsection*{Your turn}
\begin{itemize}
\raggedright
  \item \emph{Say it:} \emph{chūchāi}
  \item \emph{Say it:} \emph{chūchāi}, once more
  \item \emph{Say it:} \emph{bàba hòutiān yào chūchāi} — Dad is going away on business the day after tomorrow.
  \item \emph{Recall:} say the Mandarin for to go home from abroad, then the Mandarin for to go abroad, then say \emph{chūchāi} again
\end{itemize}

\begin{tcolorbox}[breakable,colback=teal!4,colframe=teal!35!black,title={Before you move on}]
What does chūchāi mean? (\textquotedblleft{}To travel for work; to go on a business trip\textquotedblright{}.) Say it once more.
\end{tcolorbox}
